\documentclass[11pt]{article}
\usepackage[a4paper,margin=27mm]{geometry}
\usepackage{amsmath,amssymb,amsthm,mathtools}
\usepackage{booktabs,longtable}
\usepackage[hidelinks]{hyperref}
\usepackage{microtype}

\newtheorem{theorem}{Theorem}[section]
\newtheorem{lemma}[theorem]{Lemma}
\newtheorem{proposition}[theorem]{Proposition}
\newtheorem{corollary}[theorem]{Corollary}
\theoremstyle{remark}
\newtheorem{remark}[theorem]{Remark}
\newcommand{\F}{\mathbb F_4}
\newcommand{\wt}{\operatorname{wt}}
\newcommand{\tr}{\operatorname{Tr}}
\newcommand{\im}{\operatorname{im}}
\newcommand{\Span}{\operatorname{span}_{\F}}
\newcommand{\ip}[2]{\left\langle #1,#2\right\rangle}
\newcommand{\supp}{\operatorname{supp}}
\newcommand{\Arow}[1]{\mathcal A_{#1}}
\newcommand{\Trow}[1]{\mathcal T_{#1}}
\newcommand{\Drow}[1]{\mathcal D_{#1}}

\title{Nonrealizability of a 23-qubit $M_3$-code weight enumerator}
\author{}
\date{Research draft --- 24 November 2026}

\begin{document}
\maketitle

\begin{abstract}
We give a finite nonexistence argument for the particular length-23 weight
enumerator in Appendix B, equation (B.2), of Kalra and Prakash,
\texttt{arXiv:2501.10163v3}.  If a Hermitian self-orthogonal quaternary linear
code with this enumerator existed, puncturing the support of a weight-six
word would produce a length-17 code.  Let $e$ be the number of its projective
weight-two words.  A structural argument using the support graph of these
words and the weight-ten coefficients of contained subcodes gives $e\leq4$.
An exact nonnegative identity obtained from four local fiber types and
MacWilliams identities gives
$e\geq30850111/6429420>4$.  The contradiction excludes the stated enumerator.
The argument includes a self-contained Pauli-projector positivity lemma and
an integer linear-combination certificate; its verification does not require
a numerical optimization tolerance.
\end{abstract}

\noindent\textbf{Scope and status.}
This is a research draft concerning exactly the polynomial specified below.
It does not settle the general $T$-state distillation threshold problem, nor
does it exclude every 23-qubit $M_3$ code or every other candidate enumerator.
No claim of priority, publication, or peer-review acceptance is made.

\section*{Related work and novelty scope}

The ingredients of the argument have substantial prior literature.
Residual-code methods are standard; Bouyukliev and {\"O}sterg{\aa}rd
use them in their classification of self-orthogonal codes
\cite[Section 3]{BO}.  Their Table 2 already implies the stronger
statement that no Hermitian self-orthogonal $[10,4,6]_4$ code exists.
The short-code lemma below is retained to make this proof self-contained,
not as a new small-parameter nonexistence result.
Quantum weight-enumerator linear-programming bounds have an established
history, including shadow enumerators \cite{Rains}.  Signed quantum
weight enumerators for magic-state distillation were developed by
Rall \cite{Rall}, and the positivity viewpoint used here is part of the
framework of Kalra and Prakash \cite{KP}.

The closest recent methodological comparison found in our literature
check is Lai \cite[Sections VI--VII]{Lai}: a fixed reference stabilizer,
refined MacWilliams conditions, multiplication-orbit restrictions, and
low-weight combinatorial arguments are already present there.
The length-23 candidate discussed in its Section VII.D has parameters
$[[23,1,9]]$ and vanishing coefficients through weight eight;
it differs from the present polynomial, which has $A_6=90$.
Neither residual codes, refined enumerators, nor linear programming
is claimed as a new general method in this manuscript.

The contribution claimed here is the explicit exclusion of the precise
polynomial in \cite[equation (B.2)]{KP}, by the structural bound $e\leq4$
and the displayed exact certificate forcing $e>4$.
A source check dated 24 November 2026, including the current source
version and eight distinct forward-citing papers returned by INSPIRE
and OpenAlex, found no prior resolution of this exact candidate.
This bounded search does not establish priority.  Nor do we assert
that the certificate is independent of, or stronger than, every
previously proposed refined linear or semidefinite programming bound.

\section{Statement and notation}

Write
\[
 \F=\{0,1,\omega,\omega^2\},\qquad
 \omega^2=\omega+1,\qquad \overline a=a^2.
\]
On $\F^n$ use the Hermitian form
\[
 \ip{c}{d}=\sum_{i=1}^n c_i\overline{d_i}.
\]
A code in this manuscript is an $\F$-linear subspace.  Its Hermitian dual is
$C^\perp=\{v:\ip{v}{c}=0\text{ for all }c\in C\}$.  It is
Hermitian self-orthogonal when $C\subseteq C^\perp$.  Define the ordinary
weight enumerator and its coefficients by
\[
 W_C(y)=\sum_{c\in C}y^{\wt(c)}=\sum_{j=0}^n A_j(C)y^j.
\]
Since every nonzero element of $\F$ has cube one,
$\ip{c}{c}=\wt(c)\pmod2$.  In particular, every word of a self-orthogonal
code has even weight.

The polynomial considered here is
\begin{align}
 A(y)={}&1+90y^6+1314y^8+348y^{10}+107280y^{12}
       +434880y^{14}\notag\\
       &+1282869y^{16}+1543428y^{18}
       +738072y^{20}+86022y^{22}.
 \label{eq:target}
\end{align}
It is equation (B.2) of \cite[Appendix B]{KP}, with the first homogeneous
variable specialized to one.  That paper presents it as a candidate whose
realizability is undetermined.  The candidate is relevant because its
formal distillation behavior would improve upon the five-qubit protocol.
The related quadratic-residue-code approach is developed in \cite{ZJS};
no assertion about that code family is needed in the proof below.

\begin{theorem}\label{thm:main}
There is no Hermitian self-orthogonal $\F$-linear code
$C\subseteq\F^{23}$ with weight enumerator \eqref{eq:target}.
Consequently, the 23-qubit $M_3$-code weight enumerator in
\cite[equation (B.2)]{KP} is not realizable.
\end{theorem}

Here the coding-theoretic assertion implies the $M_3$ assertion as follows.
Label the one-qubit Paulis $X,Y,Z$ by $1,\omega,\omega^2$, respectively.
Pauli multiplication modulo phase is addition in $\F$.  A transversal
Clifford cycling $X\mapsto Y\mapsto Z\mapsto X$ acts by multiplication
by $\omega$.  Therefore the additive label code of an invariant stabilizer
is $\F$-linear.  Commutativity gives
$\tr_{\F/\mathbb F_2}\ip{c}{d}=0$ for all labels $c,d$.  Applying the
same condition after multiplying one label by $\omega$ forces
$\ip{c}{d}=0$, because these are the two independent binary trace
conditions on $\F$.  A one-logical-qubit stabilizer on 23 qubits has
$2^{22}=4^{11}$ labels, so its quaternary dimension is 11.
Equivalently, the contradiction below can be read entirely as a statement
about classical quaternary codes.

\section{Two elementary coding identities}

We record the identities used in the proof, including their derivations.
This also fixes all normalization and phase conventions.

\subsection{MacWilliams identities}

For length $n$, define the quaternary Krawtchouk coefficient
\begin{equation}
 K_j^{(n)}(b)=
 \sum_{s=0}^{j}(-1)^s3^{j-s}
     \binom{b}{s}\binom{n-b}{j-s},
 \label{eq:krawtchouk}
\end{equation}
where a binomial coefficient is zero outside its usual range.  Thus
$K_j^{(n)}(b)$ is the coefficient of $y^j$ in
$(1+3y)^{n-b}(1-y)^b$.

\begin{lemma}[MacWilliams identity]\label{lem:macwilliams}
For an $[n,k]_4$ code $E$,
\begin{equation}
 W_{E^\perp}(y)=4^{-k}
 \sum_{b=0}^n A_b(E)(1+3y)^{n-b}(1-y)^b.
 \label{eq:macwilliams}
\end{equation}
In particular,
$4^k A_j(E^\perp)=\sum_b K_j^{(n)}(b)A_b(E)$.
\end{lemma}

\begin{proof}
Let $\chi(a)=(-1)^{\tr_{\F/\mathbb F_2}(a)}$.  Character orthogonality
on the linear space $E$ gives
\[
 \boldsymbol1_{E^\perp}(v)=\frac1{|E|}\sum_{c\in E}\chi(\ip{v}{c}).
\]
For completeness, if $v\notin E^\perp$, there is a $c\in E$ with
$\ip{v}{c}\ne0$; scalar closure of $E$ then supplies a nontrivial
character on $E$, whose sum vanishes.  Summing the displayed indicator
against $y^{\wt(v)}$ factors coordinatewise.  A zero coordinate of $c$
contributes $1+3y$; a nonzero coordinate contributes $1-y$, because the
sum of a nontrivial additive character over $\F^*$ is $-1$.
This proves \eqref{eq:macwilliams}.
\end{proof}

\subsection{A self-contained positivity constraint}

Let $P(c)$ be the unsigned tensor product obtained from
$0\mapsto I$, $1\mapsto X$, $\omega\mapsto Y$ and $\omega^2\mapsto Z$.

\begin{lemma}[Pauli-projector positivity]\label{lem:positivity}
Every Hermitian self-orthogonal $\F$-linear code $E$ satisfies
\begin{equation}
 W_E(i/\sqrt3)\geq0.
 \label{eq:positivity}
\end{equation}
The expression is real because every weight in $E$ is even.
\end{lemma}

\begin{proof}
For $c\in E$ set
\[
 S(c)=(-1)^{\wt(c)/2}P(c).
\]
We first check that $S(c)S(d)=S(c+d)$.  Let $s$ count coordinates at
which $c_i=d_i\ne0$.  Let $p$ and $q$ count coordinates with distinct
nonzero entries in the two cyclic orders of $X,Y,Z$, chosen so that
\[
 P(c)P(d)=i^{p-q}P(c+d).
\]
In the Hermitian product the corresponding nonzero local contributions
are $1$, $\omega$, and $\omega^2$, with the latter two possibly
interchanged according to the cyclic convention.  The equation
$\ip{c}{d}=0$ therefore says that $s,p,q$ have the same parity.
Furthermore,
\[
 \wt(c)+\wt(d)-\wt(c+d)=2s+p+q.
\]
It follows that the phase relating $S(c)S(d)$ to $S(c+d)$ is
\[
 (-1)^s i^{p+q}i^{p-q}=(-1)^{s+p}=1.
\]
Consequently $\{S(c):c\in E\}$ is a commuting Pauli subgroup and does
not contain $-I$: its only scalar label is $c=0$, for which $S(0)=I$.
Thus
\[
 \Pi_E=\frac1{|E|}\sum_{c\in E}S(c)
\]
is a Hermitian idempotent and hence a positive semidefinite projector.

Let $\lvert\tau\rangle$ be the pure qubit state with Bloch vector
$(1,1,1)/\sqrt3$.  Its expectations of $X,Y,Z$ are all $1/\sqrt3$.
Therefore
\[
 0\leq
 \langle\tau\rvert^{\otimes n}\Pi_E\lvert\tau\rangle^{\otimes n}
 =\frac1{|E|}\sum_{c\in E}(-1)^{\wt(c)/2}3^{-\wt(c)/2}
 =\frac1{|E|}W_E(i/\sqrt3).
\]
This establishes the claim for every dimension of $E$.
\end{proof}

The following consequence will exclude large components in a residual
support graph.  It is stated with exactly the hypotheses needed here, so
no bound on general code parameters is required.  Its conclusion is
already implied by \cite[Table 2]{BO}; the proof below avoids dependence
on that classification.

\begin{lemma}[Forbidden short subcode]\label{lem:short-code}
There is no Hermitian self-orthogonal code of dimension four, length at
most ten, minimum weight at least six, and at most 90 words of weight six.
\end{lemma}

\begin{proof}
Pad with zero coordinates to length ten.  Every nonzero weight is in
$\{6,8,10\}$, so write
\[
 W_E(y)=1+a y^6+b y^8+c y^{10},\qquad a+b+c=4^4-1=255.
\]
All three coefficients are nonnegative integers and $a\leq90$.
The first MacWilliams coefficient is
\[
 A_1(E^\perp)=\frac{a-c-60}{32}.
\]
It is a nonnegative integer, but it is at most $30/32<1$.
It must therefore vanish.  Hence $a=c+60$ and $b=195-2c$.
The next two dual coefficients become
\[
 A_2(E^\perp)=\frac{c-30}{4},\qquad
 A_3(E^\perp)=30-c.
\]
Their nonnegativity forces $c=30$, $a=90$, and $b=135$.
However,
\[
 W_E(i/\sqrt3)=1-\frac{90}{27}+\frac{135}{81}
                   -\frac{30}{243}=-\frac{64}{81}<0,
\]
contradicting Lemma~\ref{lem:positivity}.
\end{proof}

\section{The residual code and its Hermitian quotient}

Assume from now on that a code $C$ as in Theorem~\ref{thm:main} exists.
The target polynomial gives $|C|=4^{11}$ and minimum weight six.
Since $\dim C=11$ and $C\subseteq C^\perp\subseteq\F^{23}$,
the Hermitian quotient $C^\perp/C$ is a nondegenerate one-dimensional
space.  Every nonzero vector in such a space has nonzero Hermitian norm.
It follows that
\begin{equation}
 \{c\in C^\perp:\wt(c)\text{ is even}\}=C.
 \label{eq:even-dual}
\end{equation}
MacWilliams substitution of \eqref{eq:target} gives
\begin{align}
 W_{C^\perp}(y)-A(y)={}&927y^7+1020y^9+102600y^{11}
          +609372y^{13}\notag\\
       &+2591118y^{15}+4577364y^{17}+3738960y^{19}\notag\\
       &+924372y^{21}+37179y^{23}.
 \label{eq:dual-polynomial}
\end{align}
In particular, $C^\perp$ also has minimum weight six.
Only this minimum-weight conclusion and \eqref{eq:even-dual} will be
needed; all coefficients in \eqref{eq:dual-polynomial} follow directly
from Lemma~\ref{lem:macwilliams}.

Choose $u\in C$ of weight six, let $S=\supp(u)$, and denote the other
17 coordinates by $T$.  Coordinate permutations and nonzero coordinate
scalings preserve weights and the Hermitian form, because $a\bar a=1$
for $a\in\F^*$.  We may therefore suppose
\[
 u|_S=\boldsymbol1=(1,1,1,1,1,1),\qquad u|_T=0.
\]
Let $\pi_T$ denote puncturing the six coordinates in $S$, and set
\[
 R=\pi_T(C)\subseteq\F^{17}.
\]

\begin{lemma}\label{lem:residual}
The following statements hold.
\begin{enumerate}
\item The subcodes of $C$ and of $C^\perp$ supported in $S$ are both
      exactly $\Span\{u\}$.
\item $R$ has dimension ten, and $C|_S=\boldsymbol1^\perp\subseteq\F^6$.
\item The space
      $V=\boldsymbol1^\perp/\Span\{\boldsymbol1\}$ is nondegenerate
      Hermitian of dimension four.  There is a surjective linear map
      \[
       \phi:R\longrightarrow V,\qquad
       \phi(\pi_T(c))=[c|_S],
      \]
      satisfying $\ip{\phi(r)}{\phi(s)}=\ip{r}{s}$.
\item Its kernel $H=\ker\phi$ is the shortening of $C$ on $S$,
      viewed on $T$.  It has dimension six and minimum weight at least six.
      Moreover $R^\perp$ has dimension seven and its even-weight words
      are exactly $H$.
\end{enumerate}
\end{lemma}

\begin{proof}
Any vector supported in $S$ and not proportional to $u$ has, after
subtracting a suitable multiple of $u$, a nonzero word of weight at most
five.  Such a vector cannot belong to either $C$ or $C^\perp$, because
both codes have minimum weight six.  This proves the first assertion
and gives $\dim R=11-1=10$.

The annihilator of the restriction $C|_S$ is the subcode of $C^\perp$
supported in $S$, which is $\Span\{\boldsymbol1\}$.  Thus
$C|_S=\boldsymbol1^\perp$ and has dimension five.  The vector
$\boldsymbol1$ is isotropic, and the radical of
$\boldsymbol1^\perp$ is $\Span\{\boldsymbol1\}$.  Hence its quotient
$V$ is nondegenerate of dimension four.

Two lifts in $C$ of the same residual vector differ by a multiple of $u$,
so $\phi$ is well-defined.  The restriction result makes it surjective.
For lifts $(a,r),(b,s)\in C$,
\[
 \ip{a}{b}+\ip{r}{s}=0.
\]
Since the characteristic is two, this says
$\ip{\phi(r)}{\phi(s)}=\ip{r}{s}$.

If $\phi(r)=0$, a multiple of $u$ can be subtracted from a lift to
make its $S$ coordinates zero.  This identifies $\ker\phi$ with the
shortening of $C$.  Its dimension is $11-5=6$ and its minimum weight
is at least six.  The puncturing--shortening duality follows directly
from the definition of orthogonality:
\[
 R^\perp=\{v\in\F^{17}:(0_S,v)\in C^\perp\}.
\]
It has dimension $17-10=7$.  By \eqref{eq:even-dual}, its even-weight
words are precisely those with $(0_S,v)\in C$, namely $H$.
\end{proof}

\section{Four local fibers and 32 exact equations}

\subsection{Local classification}

For $a\in\boldsymbol1^\perp$, let $n_0,n_1,n_\omega,n_{\omega^2}$
be its four symbol multiplicities.  The condition
$\sum_i a_i=0$ says that the three nonzero multiplicities have the same
parity.  Their sum with $n_0$ is six, so all four multiplicities have
the same parity.  Up to permutation, the possibilities are
\[
 (6,0,0,0),\quad(4,2,0,0),\quad(2,2,2,0),\quad(3,1,1,1).
\]
The local fiber $a+\Span\{\boldsymbol1\}$ has weight polynomial
$\sum_{\alpha\in\F}y^{6-n_\alpha}$.  We number the four types
$0,1,2,3$, respectively, and obtain
\begin{equation}
\begin{aligned}
 f_0(y)&=1+3y^6,&
 f_1(y)&=y^2+y^4+2y^6,\\
 f_2(y)&=3y^4+y^6,&
 f_3(y)&=y^3+3y^5.
\end{aligned}
\label{eq:fibers}
\end{equation}
There are respectively
\begin{equation}
 (N_0,N_1,N_2,N_3)=(1,45,90,120)
 \label{eq:type-counts}
\end{equation}
such fibers in $V$.  Indeed, the numbers of representatives in
$\boldsymbol1^\perp$ are
\[
 4,\qquad 4\cdot3\binom62=180,\qquad
 4\frac{6!}{2!2!2!}=360,\qquad
 4\frac{6!}{3!}=480,
\]
and each fiber has four representatives.
Type zero consists only of the zero class in $V$.

For $r\in R$, its four lifts in $C$ form precisely one local fiber
and all have outside weight $\wt(r)$.  Because $C$ is even, this
outside weight is even for types $0,1,2$ and odd for type $3$.
The minimum weight of $C$ further imposes the following possibilities
for a nonzero residual vector:
\begin{equation}
\begin{aligned}
 I_0&=\{6,8,10,12,14,16\},\\
 I_1&=\{4,6,8,10,12,14,16\},\\
 I_2&=\{2,4,6,8,10,12,14,16\},\\
 I_3&=\{3,5,7,9,11,13,15,17\}.
\end{aligned}
\label{eq:allowed}
\end{equation}
In particular, $R$ has no weight-one words.  The residual vector zero
has type zero and accounts for the polynomial $f_0(y)$.

\subsection{Variables and equations}

For $b\in I_t$, let $x_{t,b}$ be the number of vectors $r\in R$ with
$\wt(r)=b$ and $\phi(r)$ of type $t$.  Put
\[
 z_{t,b}=\frac{x_{t,b}}3.
\]
Scalar multiplication by $\F^*$ preserves outside weight and local
type and acts freely on nonzero vectors, so these are nonnegative
integers.  Their integrality will not be needed for the linear certificate.
For odd $j\in\{7,9,11,13,15,17\}$ put
\[
 q_j=\frac{A_j(R^\perp)}3.
\]
These are also nonnegative integers.  There are 29 variables $z_{t,b}$
and six variables $q_j$, for a total of 35.

We now specify 32 equalities with integer coefficients.  Their labels
will be used in the certificate appendix.  Let $A_w$ denote the
coefficient of $y^w$ in \eqref{eq:target}; absent coefficients are zero.
For $w=2,4,\ldots,22$, the coefficient equations
\begin{equation}
 \Arow{w}:\qquad
 \sum_{t=0}^3\sum_{b\in I_t}
       [y^{w-b}]f_t(y)\,z_{t,b}
 =\frac{A_w-3\boldsymbol1_{w=6}}3
 \label{eq:Arows}
\end{equation}
follow from
\[
 A(y)=f_0(y)+3\sum_{t=0}^3\sum_{b\in I_t}z_{t,b}y^b f_t(y).
\]
The two equations for $w=2,4$ are included even though both are trivial.

Since $\phi$ is surjective with $4^6$ elements in every fiber, the four
type-total equations are
\begin{equation}
 \Trow{t}:\qquad
 \sum_{b\in I_t}z_{t,b}
   =\frac{4^6N_t-\boldsymbol1_{t=0}}3,
 \qquad t=0,1,2,3.
 \label{eq:Trows}
\end{equation}

Write $K_j(b)=K_j^{(17)}(b)$ and define the following linear expressions
in the variables:
\begin{equation}
 h_j=
 \begin{cases}
   z_{0,j},&j\in\{6,8,10,12,14,16\},\\
   q_j,&j\in\{7,9,11,13,15,17\},\\
   0,&j\in\{1,2,3,4,5\}.
 \end{cases}
 \label{eq:hdef}
\end{equation}
For every $1\leq j\leq17$ one has $A_j(R^\perp)=3h_j$.
For even $j$ this is Lemma~\ref{lem:residual}: the even words of
$R^\perp$ are $H$, and $r\in H$ is exactly local type zero.
For odd $j<7$ it follows from the minimum weight of $C^\perp$.
The 17 residual MacWilliams equations are therefore
\begin{equation}
 \Drow{j}:\qquad
 4^{10}h_j-
 \sum_{t=0}^3\sum_{b\in I_t} K_j(b)z_{t,b}
   =\frac{K_j(0)}3,
 \qquad 1\leq j\leq17.
 \label{eq:Drows}
\end{equation}
The factor $4^{10}$ is $|R|$.  The constant term $K_j(0)/3$ accounts
for the zero word of $R$; it is an integer for all these $j$.
Equations \eqref{eq:Arows}, \eqref{eq:Trows}, and \eqref{eq:Drows}
are the promised $11+4+17=32$ exact equalities.

Only type two can occur at residual weight two.  Consequently
\begin{equation}
 e:=\frac{A_2(R)}3=z_{2,2}
 \label{eq:e}
\end{equation}
counts the one-dimensional subspaces of $R$ generated by weight-two
words.

\section{A structural upper bound on the weight-two count}

\subsection{The support graph}

Form a graph on the 17 residual coordinate positions.  Join two positions
when they are the support of a weight-two word of $R$.

\begin{lemma}\label{lem:graph}
There is exactly one projective class of weight-two words for each edge.
Every nontrivial connected component is a complete graph.  After scaling
its coordinates, the subcode generated by the edges of a component on
$m$ vertices is
\[
 K_m=\{(a_1,\ldots,a_m)\in\F^m:a_1+\cdots+a_m=0\},
\]
of dimension $m-1$.  Subcodes from distinct components have disjoint
supports and are mutually Hermitian orthogonal.
\end{lemma}

\begin{proof}
Two independent words on the same two positions could be linearly
combined to cancel one of them, giving a weight-one word.  This is
impossible by \eqref{eq:allowed}.  Similarly, two edges sharing a
vertex can be combined to cancel that vertex and produce the edge
between their other endpoints.  This implies that every connected
component is complete.

Fix a vertex in such a component.  The words on edges from that vertex
to the other $m-1$ vertices are independent.  Scale the coordinates so
that each is a scalar multiple of the vector with ones at its two
endpoints.  They then span the displayed single-parity-check code.
Uniqueness on each edge ensures that all other edge words lie in this
same subcode.  Coordinate scalings preserve the Hermitian form, and
disjoint supports imply orthogonality.
\end{proof}

\begin{lemma}\label{lem:noK4}
Every nontrivial component of the support graph is $K_2$ or $K_3$.
\end{lemma}

\begin{proof}
A component with at least four vertices contains the three-dimensional
single-parity-check subcode generated by the edges on any four of its
vertices.  Its full inverse image in $C$ under $\pi_T$ has dimension
$3+1=4$, because $\ker\pi_T=\Span\{u\}$.  It is supported on those
four outside positions together with the six positions of $S$, so has
length at most ten.  As a subcode of $C$, it is Hermitian
self-orthogonal, has minimum weight at least six, and has at most
$A_6(C)=90$ words of weight six.  This contradicts
Lemma~\ref{lem:short-code}.
\end{proof}

\begin{proposition}\label{prop:six}
One has $e\leq6$.  If $e=5$, the nontrivial components are exactly
five copies of $K_2$.  If $e=6$, they are exactly two copies of $K_3$.
\end{proposition}

\begin{proof}
Let $a$ and $b$ be the numbers of $K_3$ and $K_2$ components.  Their
edge counts give $e=3a+b$.

The standard $K_3$ is the orthogonal complement of the all-one vector
in $\F^3$.  That vector has norm one, so $K_3$ is nondegenerate of
dimension two.  Because $\phi$ preserves the Hermitian form, its
restriction to any nondegenerate subspace is injective.  Thus the
images of the $a$ disjoint $K_3$ subcodes form a nondegenerate
$2a$-dimensional subspace of the four-dimensional space $V$, and
$a\leq2$.

Each $K_2$ is a one-dimensional zero-form space.  Its generator has a
nonzero image under $\phi$, since a zero image would be a weight-two
word of $H$, whereas $H$ has minimum weight at least six.  Images of
generators of distinct $K_2$ components are not proportional: otherwise
a nonzero linear combination would be a word of $H$ with weight four.
These images are isotropic, pairwise orthogonal, and orthogonal to all
the $K_3$ images.

If $a=2$, the nondegenerate $K_3$ image fills $V$, so $b=0$.
If $a=1$, the $K_2$ images lie in its nondegenerate
two-dimensional orthogonal complement.  Their span is totally
isotropic and has dimension at most one, so projective distinctness
gives $b\leq1$.
If $a=0$, their totally isotropic span has dimension at most two.
Indeed, in a nondegenerate Hermitian space a totally isotropic
subspace $W$ satisfies $W\subseteq W^\perp$, giving
$2\dim W\leq\dim V$.  A two-dimensional $\F$-space has only
$(4^2-1)/(4-1)=5$ projective points, so $b\leq5$.
These alternatives prove the assertion and the classifications at
$e=5$ and $e=6$.
\end{proof}

\subsection{Excluding five disjoint edges}

\begin{proposition}\label{prop:no-five}
The case $e=5$ is impossible.
\end{proposition}

\begin{proof}
Suppose the support graph has five $K_2$ components.  Let $K\subseteq R$
be the subcode they generate.  After coordinate scalings, $K$ is the
pair-repetition image of $\F^5$: each of five symbols is repeated on
its own pair of coordinates.  It has dimension five, support ten, and
identically zero Hermitian form.

Put $W=\phi(K)$.  By the proof of Proposition~\ref{prop:six}, the five
pair generators have nonzero, projectively distinct images spanning a
totally isotropic space of dimension at most two.  They must therefore
span a two-dimensional space and represent all five projective points
of $W$.  Each generator has residual weight two, so its local type is
two.  Local type is invariant under nonzero scalar multiplication.
It follows that every nonzero vector of $W$ has local type two, while
zero has type zero.

In the five pair coordinates, the kernel of $\phi|_K$ is a
three-dimensional code $L\subseteq\F^5$ whose parity-check columns
represent the five distinct projective points of $\F^2$.  Any nonzero
linear functional on $\F^2$ vanishes on exactly one of these five
points.  Hence the corresponding two-dimensional check code has
enumerator $1+15z^4$, and MacWilliams gives
\begin{align}
 h(z):=W_L(z)
 &=\frac{(1+3z)^5+15(1+3z)(1-z)^4}{16}\notag\\
 &=1+30z^3+15z^4+18z^5.
 \label{eq:hamming}
\end{align}

Let $U=\{c\in C:\pi_T(c)\in K\}$.  It has dimension six and
support at most 16.  A vector in the five pair coordinates with
weight $j$ has outside weight $2j$.  If it belongs to $L$, its local
fiber is $f_0$; otherwise its local fiber is $f_2$.  Thus
\begin{equation}
 W_U(y)=h(y^2)(1+3y^6)
       +\bigl((1+3y^2)^5-h(y^2)\bigr)(3y^4+y^6).
 \label{eq:e5-enumerator}
\end{equation}
The coefficient at weight ten is
\[
 A_{10}(U)=18+\binom52 3^2
                +3\left(\binom53 3^3-30\right)
          =18+90+720=828.
\]
The three summands count, respectively, kernel words of pair weight
five with local weight zero, nonkernel words of pair weight two with
local weight six, and nonkernel words of pair weight three with local
weight four.  Since $U\subseteq C$ but $A_{10}(C)=348$, this is a
contradiction.
\end{proof}

\subsection{Excluding two triangles}

\begin{proposition}\label{prop:no-six}
The case $e=6$ is impossible.  Consequently, every code satisfying
the assumed target enumerator would have $e\leq4$.
\end{proposition}

\begin{proof}
Suppose the support graph consists of two nontrivial $K_3$ components.
Their generated subcode is
$K=K_3\oplus K_3\subseteq R$, of dimension four.  It is nondegenerate,
so the form-preserving map $\phi|_K$ is an isomorphism onto $V$.
The enumerator of one $K_3$ is $1+9y^2+6y^3$.  Therefore
\[
 W_K(y)=(1+9y^2+6y^3)^2
 =1+18y^2+12y^3+81y^4+108y^5+36y^6.
\]
Residual weight-two vectors have local type two; vectors of outside
weight three or five have local type three.  At outside weights four
and six, the local type is one or two.  It cannot be zero, because
$\phi|_K$ is injective.  Since $\phi(K)=V$, the total counts of local
types are exactly $(1,45,90,120)$.

Let $a$ be the number of vectors of $K$ with outside weight four and
local type one.  The remaining even-weight type counts are forced:
\[
\begin{array}{c|cc}
 &\text{type one}&\text{type two}\\\hline
 \text{outside weight four}&a&81-a\\
 \text{outside weight six}&45-a&a-9
\end{array}
\]
Thus $9\leq a\leq45$.

The full inverse image $U=\{c\in C:\pi_T(c)\in K\}$ has dimension
five and support at most 12.  Using the fiber polynomials in
\eqref{eq:fibers}, its weight-ten coefficient is
\begin{align*}
 A_{10}(U)
 &=2a+(81-a)+3\cdot108+(45-a)+3(a-9)\\
 &=423+3a\geq450.
\end{align*}
The five summands correspond to outside/local types
$(4,1),(4,2),(5,3),(6,1),(6,2)$, respectively.
This exceeds $A_{10}(C)=348$, again a contradiction.
Propositions~\ref{prop:six} and \ref{prop:no-five} now yield $e\leq4$.
\end{proof}

\section{The exact linear certificate and the contradiction}

\begin{lemma}\label{lem:certificate}
Every solution of the 32 equalities
\eqref{eq:Arows}, \eqref{eq:Trows}, and \eqref{eq:Drows} satisfies
\begin{align}
 12858840\,e={}&61700222
       +415155z_{0,10}+2259z_{0,12}+83404z_{0,16}\notag\\
       &+1207628z_{1,4}+333082z_{1,8}+890442z_{2,6}\notag\\
       &+698464z_{3,5}+251078z_{3,7}+187648q_9.
 \label{eq:certificate}
\end{align}
\end{lemma}

\begin{proof}
This is an exact linear combination of the displayed integer
equalities.  Appendix~\ref{app:certificate} gives all nonzero integer
multipliers and an explicit normalization, so that the identity can
be verified by coefficient collection using
\eqref{eq:krawtchouk}.  No inequalities, coefficient upper bounds,
or integrality assumptions enter this calculation.
\end{proof}

\begin{proof}[Proof of Theorem~\ref{thm:main}]
If the target code existed, all $z_{t,b}$ and $q_j$ in
\eqref{eq:certificate} would be nonnegative.  Thus
\[
 e\geq\frac{61700222}{12858840}
   =\frac{30850111}{6429420}>4.
\]
On the other hand, Proposition~\ref{prop:no-six} gives $e\leq4$.
These two necessary conditions are incompatible.  This proves
nonexistence of the code and hence of the stated $M_3$-code enumerator.
\end{proof}

\begin{remark}[What the certificate uses]
The inequality $e>4$ follows already over the real nonnegative cone.
In particular, there is no need to round a numerical optimum, solve an
integer program, or assume any finite upper bounds on residual
coefficients.  The exact excess over four is
\[
 \frac{30850111}{6429420}-4=\frac{5132431}{6429420}>0.
\]
The structural argument is essential to this proof: the four-type linear constraints
alone do not encode the forbidden subcodes responsible for $e\leq4$.
\end{remark}

\section{Interpretation and limitations}

The theorem removes one specific formal enumerator from the space of
possible $T$-state distillation codes.  Its mechanism is local: a
weight-six stabilizer word forces a four-dimensional Hermitian quotient,
and the low-weight geometry of the residual code controls subcodes
whose weight-ten counts are too large.  The linear certificate then
shows that the enumerator necessarily demands more weight-two
structure than those subcode restrictions permit.

\paragraph{Formal verification of the finite certificate.}
The accompanying Lean file \texttt{N23Certificate.lean} has been checked
locally with Lean 4.19.0 (compiler commit \texttt{6caaee842e94}).
It verifies the fixed $32\times35$ integer coefficient identities and
proves that the encoded row equations, nonnegativity of the 35 variables,
and $e\leq4$ are inconsistent.  There are no admitted goals or
user-declared axioms; the final implication uses only Lean's standard
logical foundations.  The derivation of these equations from a
quaternary code, the Hermitian quotient geometry, and the proof of
$e\leq4$ remain mathematical arguments in this manuscript.
Thus the formal check covers the finite arithmetic implication,
while the coding-theoretic reductions retain their written proofs.

Nothing in the proof supplies a universal upper bound on the
$T$-direction distillation threshold or rules out a different enumerator
at length 23 or at another length.  Publication-level novelty and the
relationship to any independently obtained results require separate
literature and expert review.  The mathematical claim of this draft is
exactly Theorem~\ref{thm:main}.

\appendix
\section{Integer multipliers for the certificate}\label{app:certificate}

Order the 35 variables as $z_{t,b}$ with $t=0,1,2,3$ and increasing
$b\in I_t$, followed by $q_7,q_9,q_{11},q_{13},q_{15},q_{17}$.
Let $v$ be this column vector.  Let $Mv=d$ be the 32-row system whose
rows, in order, are $\Arow{2},\Arow{4},\ldots,\Arow{22}$,
$\Trow{0},\ldots,\Trow{3}$, and
$\Drow{1},\ldots,\Drow{17}$, with the signs exactly as written in
\eqref{eq:Arows}--\eqref{eq:Drows}.
All its matrix entries and right-hand sides are integers.

Assign the following integer multipliers $\lambda$; every unlisted row
has multiplier zero.
\begin{center}
\begin{longtable}{lr}
\toprule
Row&Multiplier\\
\midrule
\endfirsthead
\toprule
Row&Multiplier\\
\midrule
\endhead
$\Arow{6}$&$-431289105217$\\
$\Arow{8}$&$245979135951$\\
$\Arow{10}$&$-342176944369$\\
$\Arow{12}$&$4228502143$\\
$\Arow{14}$&$2237573663$\\
$\Arow{16}$&$1112061423$\\
$\Arow{18}$&$851965423$\\
$\Arow{20}$&$1457285663$\\
$\Arow{22}$&$2928022143$\\
$\Trow{0}$&$-99690634384$\\
$\Trow{1}$&$-10504369200$\\
$\Trow{2}$&$-9638952960$\\
$\Drow{3}$&$-18731648$\\
$\Drow{4}$&$-4474665$\\
$\Drow{5}$&$-238208$\\
$\Drow{6}$&$400590$\\
$\Drow{8}$&$-140874$\\
$\Drow{9}$&$-187648$\\
$\Drow{10}$&$1331$\\
$\Drow{12}$&$81145$\\
$\Drow{14}$&$83404$\\
\bottomrule
\end{longtable}
\end{center}

Let $e_{2,2}$ denote the coordinate vector selecting $z_{2,2}$, and let
$r$ be zero except at the nine variable positions listed below:
\begin{center}
\begin{tabular}{lr}
\toprule
Variable&Nonnegative coefficient in $r$\\
\midrule
$z_{0,10}$&$435321569280$\\
$z_{0,12}$&$2368733184$\\
$z_{0,16}$&$87455432704$\\
$z_{1,4}$&$1266289737728$\\
$z_{1,8}$&$349261791232$\\
$z_{2,6}$&$933696110592$\\
$z_{3,5}$&$732392587264$\\
$z_{3,7}$&$263274364928$\\
$q_9$&$196763189248$\\
\bottomrule
\end{tabular}
\end{center}
Direct integer arithmetic gives
\begin{align}
 M^{\mathsf T}\lambda
   &=13483471011840\,e_{2,2}-r,
 \label{eq:matrix-cert}\\
 d^{\mathsf T}\lambda&=64697371983872.
 \label{eq:rhs-cert}
\end{align}
Applying these identities to $Mv=d$ gives
\[
 13483471011840\,e=64697371983872+r^{\mathsf T}v.
\]
Every coefficient of this last identity is divisible by $2^{20}=1048576$,
yielding exactly \eqref{eq:certificate}.  Equations
\eqref{eq:matrix-cert} and \eqref{eq:rhs-cert}, together with
the nonnegativity of $r$, are a finite integer certificate.
Their strict incompatibility with $e\leq4$ is also visible without
division:
\[
 64697371983872>4\cdot13483471011840=53933884047360.
\]
The equality multipliers, the nine residual coefficients, and the
defining row formulas suffice to reproduce the verification without
using an optimization solver.

\begin{thebibliography}{9}

\bibitem{KP}
Amolak Ratan Kalra and Shiroman Prakash,
\emph{Invariant Theory, Magic State Distillation, and Bounds on Classical
Codes}, arXiv:2501.10163v3, revised 23 March 2026.
\url{https://arxiv.org/abs/2501.10163v3}.

\bibitem{ZJS}
Michael Zurel, Santanil Jana, and Nadish de Silva,
\emph{High-threshold magic state distillation with quantum quadratic
residue codes}, arXiv:2603.18560v2, revised 30 September 2026;
\emph{Quantum Science and Technology} \textbf{11}, 045052 (2026).
\url{https://arxiv.org/abs/2603.18560v2}.

\bibitem{BO}
Iliya Bouyukliev and Patric R. J. {\"O}sterg{\aa}rd,
\emph{Classification of Self-Orthogonal Codes over $F_3$ and $F_4$},
\emph{SIAM Journal on Discrete Mathematics} \textbf{19}(2),
363--370 (2005).
\url{https://doi.org/10.1137/S0895480104441085}.

\bibitem{Rains}
E. M. Rains,
\emph{Quantum shadow enumerators},
\emph{IEEE Transactions on Information Theory} \textbf{45}(7),
2361--2366 (1999); arXiv:quant-ph/9611001v2.
\url{https://arxiv.org/abs/quant-ph/9611001}.

\bibitem{Rall}
Patrick Rall,
\emph{Signed quantum weight enumerators characterize qubit magic state
distillation}, arXiv:1702.06990v1 (2017).
\url{https://arxiv.org/abs/1702.06990}.

\bibitem{Lai}
Ching-Yi Lai,
\emph{Twisted Fourier Analysis of Quantum Weight Enumerators},
arXiv:2610.37894v1, 27 October 2026.
\url{https://arxiv.org/abs/2610.37894}.

\end{thebibliography}
\end{document}
